% !TeX program = pdflatex
\documentclass[aspectratio=169,xcolor={dvipsnames,table}]{beamer}

\input{../../_en_preambulo_beamer}
% Comandos y macros estadísticas compartidas para las presentaciones Beamer en inglés (Metropolis)

% Conjuntos numéricos básicos
\newcommand{\R}{\mathbb{R}}
\newcommand{\N}{\mathbb{N}}
\newcommand{\Q}{\mathbb{Q}}
\newcommand{\Z}{\mathbb{Z}}
\newcommand{\set}[1]{\left\{ #1 \right\}}
\newcommand{\abs}[1]{\left|#1\right|}
\newcommand{\norm}[1]{\left\| #1 \right\|}

% Comandos estadísticos y probabilísticos del curso
\newcommand{\Var}{\operatorname{Var}}
\newcommand{\cov}{\operatorname{Cov}}
\newcommand{\corr}{\rho}
\newcommand{\comb}[2]{\begin{pmatrix} #1 \\ #2 \end{pmatrix}}
\newcommand{\s}{\sigma}
\newcommand{\particion}{\mathfrak{P}}
\newcommand{\card}[1]{\#\left( #1 \right)}
\newcommand{\seq}[3]{\set{#1}_{#2}^{#3}}
\newcommand{\E}{\mathbb{E}}
\newcommand{\Prob}{\mathbb{P}}

% Entornos pedagógicos y cajas en inglés académico natural (soportando tanto nombres en inglés como en español)
\newenvironment{discussion}{%
  \begin{alertblock}{Discussion Question}%
}{%
  \end{alertblock}%
}
\newenvironment{discusion}{%
  \begin{alertblock}{Discussion Question}%
}{%
  \end{alertblock}%
}

\newenvironment{reflection}{%
  \begin{alertblock}{Classroom Reflection}%
}{%
  \end{alertblock}%
}
\newenvironment{reflexion}{%
  \begin{alertblock}{Classroom Reflection}%
}{%
  \end{alertblock}%
}

\newenvironment{paradox}{%
  \begin{alertblock}{Exploratory Paradox}%
}{%
  \end{alertblock}%
}
\newenvironment{paradoja}{%
  \begin{alertblock}{Exploratory Paradox}%
}{%
  \end{alertblock}%
}

\newenvironment{motivation}{%
  \begin{exampleblock}{Motivation: Data Science in Practice}%
}{%
  \end{exampleblock}%
}
\newenvironment{motivacion}{%
  \begin{exampleblock}{Motivation: Data Science in Practice}%
}{%
  \end{exampleblock}%
}

\newenvironment{quickchallenge}{%
  \begin{block}{Quick Team Challenge}%
}{%
  \end{block}%
}
\newenvironment{retoclase}{%
  \begin{block}{Quick Team Challenge}%
}{%
  \end{block}%
}


\title{Multinomial Distribution}
\subtitle{Section 03.08 --- Polytomous Trials, Negative Covariances and Contingency Tables}
\author[J. Castillo Colmenares]{Juliho Castillo Colmenares}
\institute[Tec de Monterrey]{Tecnologico de Monterrey}
\date{\vspace{-1.5cm}}

\begin{document}

% SLIDE 01: Title Page
\begin{frame}[plain]
	\titlepage
\end{frame}

% SLIDE 02: Roadmap
\begin{frame}{Roadmap --- Chapter 03: Discrete Random Variables}
	\vspace{-0.15cm}
	\begin{block}{Curricular Structure of Chapter 03}
		\footnotesize
		\begin{enumerate}\itemsep=0.03cm
			\item \textcolor{gray}{Section 03.01: Discrete Probability Mass Functions (PMF and Support)}
			\item \textcolor{gray}{Section 03.02: Cumulative Distribution Functions (Discrete CDF)}
			\item \textcolor{gray}{Section 03.03: Mathematical Expectation, Variance, and Moments}
			\item \textcolor{gray}{Section 03.04: Bernoulli and Binomial Distributions}
			\item \textcolor{gray}{Section 03.05: Geometric and Negative Binomial Distributions}
			\item \textcolor{gray}{Section 03.06: Hypergeometric Distribution}
			\item \textcolor{gray}{Section 03.07: Poisson Distribution}
			\item \textbf{\textcolor{TecRojo}{Section 03.08: Multinomial Distribution}} $\leftarrow$ \emph{Current focus}
		\end{enumerate}
	\end{block}
	\vspace{-0.2cm}
	\begin{alertblock}{Learning Objective}
		\scriptsize
		Extend the Binomial distribution to trials with $k \ge 3$ mutually exclusive categories, exploit the negative covariance structure induced by the constraint $\sum_{i} X_{i} = n$, and apply the conditional Multinomial distribution to contingency table analysis.
	\end{alertblock}
\end{frame}

% SLIDE 03: Motivation and Formal Definition
\begin{frame}{From Binomial to Multinomial: Polytomous Trials}
	\vspace{-0.15cm}
	\begin{columns}[T]
		\begin{column}{0.48\textwidth}
			\begin{block}{Formal Definition}
				\footnotesize
				A random vector $\bm{X} = (X_{1}, \dots, X_{k})$ follows a Multinomial distribution $\text{Mult}(n, p_{1}, \dots, p_{k})$ if:
				$$P(\bm{X} = \bm{x}) = \dfrac{n!}{x_{1}! \cdots x_{k}!} p_{1}^{x_{1}} \cdots p_{k}^{x_{k}}$$
				with $\sum_{i=1}^{k} x_{i} = n$ and $\sum_{i=1}^{k} p_{i} = 1$.
			\end{block}
		\end{column}
		\begin{column}{0.48\textwidth}
			\begin{alertblock}{Marginal Binomial Theorem}
				\footnotesize
				Each component is marginally Binomial: $X_{i} \sim \text{Bin}(n, p_{i})$. The Multinomial generalizes the Binomial by accounting for joint probabilities across multiple categories.
			\end{alertblock}
		\end{column}
	\end{columns}
\end{frame}

% SLIDE 03B: Marginal Binomial Reduction
\begin{frame}{Marginal Reduction: Each $X_{i} \sim \text{Bin}(n, p_{i})$}
	\vspace{-0.15cm}
	\begin{block}{Marginalization Theorem}
		\footnotesize
		The Multinomial distribution is univariately Binomial: each marginal component is Binomial. The proof proceeds by summing over the remaining components.
		$$P(X_{i} = x_{i}) = \binom{n}{x_{i}} p_{i}^{x_{i}} (1 - p_{i})^{n - x_{i}}, \quad X_{i} \sim \text{Bin}(n, p_{i})$$
		The proof uses the Binomial Theorem when summing over all configurations of the other $k-1$ components.
	\end{block}
	\vspace{-0.1cm}
	\pause
	\begin{alertblock}{Marginal Expectation and Variance}
		\footnotesize
		\begin{itemize}\itemsep=0.03cm
			\item $\E[X_{i}] = n p_{i}$
			\item $\Var(X_{i}) = n p_{i} (1 - p_{i})$
			\item Marginals \emph{do not capture} the dependencies; the covariance matrix is required for joint analysis.
		\end{itemize}
	\end{alertblock}
\end{frame}

% SLIDE 04: Covariance Structure
\begin{frame}{Negative Covariance Structure: $\cov(X_{i}, X_{j}) = -n p_{i} p_{j}$}
	\vspace{-0.15cm}
	\begin{columns}[T]
		\begin{column}{0.48\textwidth}
			\begin{block}{Derivation via Indicators}
				\footnotesize
				Represent $X_{i} = \sum_{r=1}^{n} I_{r}^{(i)}$ with $I_{r}^{(i)} = 1$ if trial $r$ yields category $i$. By bilinearity and independence between trials:
				$$\cov(X_{i}, X_{j}) = n \cdot (-p_{i} p_{j}) = -n p_{i} p_{j}$$
			\end{block}
		\end{column}
		\begin{column}{0.48\textwidth}
			\begin{alertblock}{Singular Covariance Matrix}
				\footnotesize
				The covariance matrix $\bm{\Sigma}$ has diagonal $np_{i}(1-p_{i})$ and off-diagonal $-np_{i}p_{j}$. It is \emph{singular} (rank $k-1$) because $\sum_{i} X_{i} = n$ imposes a linear constraint.
			\end{alertblock}
		\end{column}
	\end{columns}
\end{frame}

% SLIDE 05: Conditional Multinomial
\begin{frame}{Conditional Multinomial Distribution}
	\vspace{-0.15cm}
	\begin{block}{Theorem: Sub-vector given Marginal}
		\scriptsize
		If $\bm{X} = (X_{1}, \dots, X_{k}) \sim \text{Mult}(n, p_{1}, \dots, p_{k})$ and $X_{k} = m_{k}$ is fixed, then:
		$$(X_{1}, \dots, X_{k-1}) \mid (X_{k} = m_{k}) \sim \text{Mult}\left(n - m_{k}, \dfrac{p_{1}}{1 - p_{k}}, \dots, \dfrac{p_{k-1}}{1 - p_{k}}\right)$$
	\end{block}
	\vspace{0.05cm}
	\pause
	\begin{alertblock}{Special Case: Conditional Binomial}
		\scriptsize
		For $k = 3$, conditioning on $X_{2} + X_{3} = m$ produces:
		$$X_{2} \mid (X_{2} + X_{3} = m) \sim \text{Bin}\left(m, \dfrac{p_{2}}{p_{2} + p_{3}}\right)$$
		This is the foundation of Fisher's exact test in contingency tables.
	\end{alertblock}
\end{frame}

% SLIDE 06: Applications
\begin{frame}{Applications in Data Science and Multivariate Analysis}
	\vspace{-0.1cm}
	\begin{itemize}\itemsep=0.1cm
		\item \textbf{Opinion Surveys and Political Science:} Distribution of preferences across multiple candidates in a sample of voters. \pause
		\item \textbf{Industrial Quality Management:} Classification of manufactured items into multiple quality categories (Premium, Standard, Defective). \pause
		\item \textbf{Population Genomics:} Genotype distribution under Hardy-Weinberg equilibrium with three genotypes (homozygous dominant, heterozygous, homozygous recessive). \pause
		\item \textbf{Natural Language Processing:} Frequency distribution of words or grammatical categories in a text corpus.
	\end{itemize}
\end{frame}

% SLIDE 07: Python Lab Block 1
\begin{frame}[fragile]{Python Lab: PMF and Marginalization}
	\vspace{-0.05cm}
	\scriptsize
	Computational verification in SciPy (\texttt{03.08\_multinomial\_distribution.py}) of the Multinomial PMF, normalization, and marginal Binomial reduction:
	\vspace{0.02cm}
	{\lstset{basicstyle=\fontsize{5pt}{6pt}\selectfont\ttfamily}
	\lstinputlisting[firstline=15, lastline=39]{../../code/03_variables_aleatorias_discretas/03.08_multinomial_distribution.py}}
\end{frame}

% SLIDE 08: Python Lab Block 2
\begin{frame}[fragile]{Python Lab: Negative Covariances}
	\vspace{-0.05cm}
	\scriptsize
	Monte Carlo simulation with $250{,}000$ trials verifying $\cov(X_{i}, X_{j}) = -n p_{i} p_{j}$:
	\vspace{0.02cm}
	{\lstset{basicstyle=\fontsize{4.5pt}{5.5pt}\selectfont\ttfamily}
	\lstinputlisting[firstline=43, lastline=68]{../../code/03_variables_aleatorias_discretas/03.08_multinomial_distribution.py}}
\end{frame}

% SLIDE 09: Python Lab Block 3
\begin{frame}[fragile]{Python Lab: Conditional Multinomial}
	\vspace{-0.05cm}
	\scriptsize
	Numerical verification that $X_{3} \mid (X_{3} + X_{4} = m) \sim \text{Bin}(m, p_{3}/(p_{3} + p_{4}))$:
	\vspace{0.02cm}
	{\lstset{basicstyle=\fontsize{4.5pt}{5.5pt}\selectfont\ttfamily}
	\lstinputlisting[firstline=83, lastline=108]{../../code/03_variables_aleatorias_discretas/03.08_multinomial_distribution.py}}
\end{frame}

% SLIDE 10: Python Lab Terminal Output
\begin{frame}[fragile]{Python Lab: Terminal Output}
	\vspace{-0.1cm}
	\begin{block}{}
		\fontsize{4pt}{5pt}\selectfont
		\begin{verbatim}
=== Block 1: Multinomial PMF & Marginalization ===
n=100, p=[0.45, 0.35, 0.20], target=(50, 30, 20)
P(X=(50,30,20)) = 4.79e-03
PMF normalization: 1.000000
Marginal X_1 ~ Bin(100, 0.45): P(X_1=50) theo=0.048152 | joint=0.048152

=== Block 2: Covariance Structure & Negative Correlations ===
Theoretical: E[X]=(50,30,20) | Var=(25,21,16)
  Cov(X_1,X_2)=-15 | Cov(X_1,X_3)=-10 | Cov(X_2,X_3)=-6
Empirical (N=250,000): E[X]=(50.0,30.0,20.0)
  Cov(X_1,X_2)=-14.96 | Cov(X_1,X_3)=-10.12 | Cov(X_2,X_3)=-5.93

=== Block 3: Conditional Distributions ===
Mult(100, [0.4, 0.3, 0.2, 0.1]) | p_cond = 0.6667
Conditional X_3 | (X_3+X_4=30) ~ Bin(30, 0.6667)
Theoretical Bin(30, 0.6667):
  k=18: 0.1101 | k=19: 0.1391 | k=20: 0.1530 | k=21: 0.1457
Empirical conditional (N=500,000, |filter|=43,324):
  k=18: emp=0.1109 | theo=0.1101 | k=20: emp=0.1497 | theo=0.1530
Problem 3.8.9 verification: P(X_3=18 | X_3+X_4=30) = 0.1101
		\end{verbatim}
	\end{block}
\end{frame}









% SLIDE 19: Closing and Next Steps
\begin{frame}{Closing Section and Modular Perspective}
	\vspace{-0.15cm}
	\begin{block}{Synthesis on Multivariate Categorical Distributions}
		\scriptsize
		The Multinomial distribution teaches us that extending the Binomial to multiple categories preserves the marginals (each $X_{i} \sim \text{Bin}(n, p_{i})$) but introduces new structured dependencies: the negative covariances $\cov(X_{i}, X_{j}) = -n p_{i} p_{j}$ encode the fundamental constraint $\sum X_{i} = n$.
	\end{block}
	\vspace{0.05cm}
	\pause
	\begin{alertblock}{Key Lessons and Next Step}
		\begin{itemize}\itemsep=0.05cm
			\item \textbf{Preserved Marginals:} Each $X_{i} \sim \text{Bin}(n, p_{i})$ individually, but joint dependencies are essential.
			\item \textbf{Negative Covariances:} Reflect the $\sum X_{i} = n$ constraint; excess in $i$ implies deficit in some $j$.
			\item \textbf{Conditional Multinomial:} Under fixed marginals, distributions collapse to Multinomials with fewer categories.
			\item \textbf{Next Session:} \textcolor{TecRojo}{Normal Distribution and Continuous Approximation of Discrete Variables}.
		\end{itemize}
	\end{alertblock}
\end{frame}

\end{document}
